\chapter{1977年广西百色卷（文）}

\section{解答题}

\begin{problem}
计算：\(-3.5-\left[2\frac{3}{7}-\left(-3\frac{3}{14}-3\frac{5}{14}\right)\right]\).
\end{problem}

\begin{solution}
先计算括号内的两个带分数，得
\(-3\frac{3}{14}-3\frac{5}{14}=-6\frac{4}{7}\)，从而
\(2\frac{3}{7}-(-6\frac{4}{7})=9\). 因此
\[
-3.5-\left[2\frac{3}{7}-\left(-3\frac{3}{14}-3\frac{5}{14}\right)\right]
=-3.5-\left(2\frac{3}{7}+6\frac{4}{7}\right)
=-3.5-9=-12.5
\]
\end{solution}

\begin{problem}
当 \(|m|=|n|\) 时，能肯定 \(m=n\) 吗？为什么？
\end{problem}

\begin{solution}
不能. 由 \(|m|=|n|\) 可得 \(m^2=n^2\)，从而
\((m-n)(m+n)=0\)，所以可能有 \(m=n\)，也可能有 \(m=-n\). 例如取
\(m=1\)，\(n=-1\)，则 \(|m|=|n|=1\)，但 \(m\ne n\). 因此仅由
\(|m|=|n|\) 不能肯定 \(m=n\).
\end{solution}

\begin{problem}
当 \(a=1\frac{1}{2}\) 时，求 \(a+\frac{1-2a+a^2}{a-1}\) 的值.
\end{problem}

\begin{solution}
由题设 \(a=1\frac{1}{2}=\frac{3}{2}\)，所以 \(a-1\ne0\)，原式中的分式有意义. 又
\[
a+\frac{1-2a+a^2}{a-1}
=a+\frac{(a-1)^2}{a-1}
=a+(a-1)=2a-1
\]
代入 \(a=\frac{3}{2}\)，得 \(2a-1=2\times\frac{3}{2}-1=2\).
\end{solution}

\begin{problem}
已知 \(\log_{x}32=5\). 求 \(x\).
\end{problem}

\begin{solution}
对数的底数必须满足 \(x>0\) 且 \(x\ne1\). 由对数定义，
\(\log_x32=5\) 等价于 \(x^5=32=2^5\). 实数范围内五次方函数严格递增，
所以 \(x=2\)，且该值满足对数底数的条件.
\end{solution}

\begin{problem}
分解下列各式的因式：

\begin{enumerate}
\item \((x-2)(x^2-3)+(x-2)^2-(x-2)(2x-3)\)；
\item \(bc^2+a^2c+ab^2-a^2b-ac^2-b^2c\).
\end{enumerate}
\end{problem}

\begin{solution}
\begin{enumerate}
\item
原式中各项都有公因式 \(x-2\)，因此
\[
(x-2)(x^2-3)+(x-2)^2-(x-2)(2x-3)
=(x-2)\left[(x^2-3)+(x-2)-(2x-3)\right]
=(x-2)(x^2-x-2)
\]
\[
=(x-2)(x-2)(x+1)=(x-2)^2(x+1)
\]

\item
设原式为 \(E\). 按含有 \(c^2\)、\(c\) 和常数项分组，得
\[
E=c^2(b-a)+c(a^2-b^2)+ab(b-a)
=(b-a)\left[c^2-c(a+b)+ab\right]
\]
括号内的二次式可分解为
\(c^2-c(a+b)+ab=(c-a)(c-b)\)，所以
\[
E=(b-a)(c-a)(c-b)
\]
\end{enumerate}
\end{solution}

\begin{problem}
已知 \(AD\) 和 \(BE\) 分别是 \(\triangle ABC\) 的 \(BC\) 和 \(AC\) 边上的高，延长 \(AD\) 交外接圆于 \(F\)，求证 \(\angle EBC=\angle FBC\).

\FigureLayoutDeclare{img/1977/guangxi_baise_liberal/q6_fig1.png}{7.6cm}{35bp 0bp 0bp 3bp}
\bitmapfigure[max width=0.34\linewidth]{img/1977/guangxi_baise_liberal/q6_fig1.png}
\end{problem}

\begin{solution}
在 \(\triangle BEC\) 和 \(\triangle ADC\) 中，\(BE\perp AC\)，\(AD\perp BC\)，所以
\(\angle BEC=\angle ADC=90^\circ\). 又因为 \(E\) 在 \(AC\) 上、\(D\) 在 \(BC\) 上，
所以 \(\angle BCE=\angle ACD\). 因此
\(\triangle BEC\sim\triangle ADC\)，对应关系为 \(B\leftrightarrow A\)，
\(E\leftrightarrow D\)，\(C\leftrightarrow C\)，从而
\(\angle EBC=\angle DAC\).

由构造知 \(A\)，\(D\)，\(F\) 三点共线，且射线 \(AD\) 与射线 \(AF\) 同向，所以
\(\angle DAC=\angle FAC\). 又因为 \(A\)，\(B\)，\(C\)，\(F\) 在同一圆上，
\(\angle FAC\) 与 \(\angle FBC\) 都是所对弦 \(FC\) 所对的圆周角，故
\(\angle FAC=\angle FBC\). 因而 \(\angle EBC=\angle FBC\).
\end{solution}

\begin{problem}
红星农场一块地，用一台拖拉机来耕 \(6\text{ 天}\)，完成一半，后来加入了另一台新式拖拉机，两台合耕 \(2\text{ 天}\) 就耕完余下的地，新式拖拉机单独耕这块地要几天？
\end{problem}

\begin{solution}
设整块地的工作量为 \(1\). 旧式拖拉机 \(6\text{ 天}\) 完成了 \(\frac{1}{2}\)，
所以它每天完成的工作量为 \(\frac{\frac{1}{2}}{6}=\frac{1}{12}\). 设新式拖拉机单独完成整块地需要
\(x\text{ 天}\)，则它每天完成的工作量为 \(\frac{1}{x}\)，其中 \(x>0\).
两台拖拉机合耕 \(2\text{ 天}\) 完成余下的 \(\frac{1}{2}\)，故
\[
2\left(\frac{1}{12}+\frac{1}{x}\right)=\frac{1}{2}
\]
解得 \(\frac{1}{12}+\frac{1}{x}=\frac{1}{4}\)，所以 \(\frac{1}{x}=\frac{1}{6}\)，
从而 \(x=6\). 因此新式拖拉机单独耕完这块地需要 \(6\text{ 天}\).
\end{solution}

\begin{problem}
曙光人民公社养鸡场，准备在鸡舍后面开辟一个矩形养鸡场，现仓库里有 \(38\text{ 米}\) 长的铁丝网，问鸡场的长和宽各是多少米？才能使鸡场的面积最大？（靠鸡舍一边不围网）
\end{problem}

\begin{solution}
设鸡场的宽为 \(x\text{ 米}\). 由于靠鸡舍的一边不围网，需要围成两条宽和一条长，
所以鸡场的长为 \(38-2x\text{ 米}\). 长、宽都为正数，故 \(0<x<19\).
设鸡场面积为 \(S\text{ 平方米}\)，则
\[
S=x(38-2x)=-2x^2+38x=-2\left(x-\frac{19}{2}\right)^2+\frac{361}{2}
\]
由于平方项不小于零，\(S\leqslant\frac{361}{2}\)，且当
\(x=\frac{19}{2}=9.5\) 时等号成立. 此时鸡场的长为
\(38-2\times9.5=19\text{ 米}\). 因此鸡场的宽为 \(9.5\text{ 米}\)、长为
\(19\text{ 米}\) 时面积最大.
\end{solution}

\begin{problem}
已知一个动点 \(P(x,y)\) 到两定点 \(A(1,3)\) 和 \(B(-3,5)\) 的距离相等，求这动点的轨迹方程.
\end{problem}

\begin{solution}
由距离相等，得
\[
\sqrt{(x-1)^2+(y-3)^2}=\sqrt{(x+3)^2+(y-5)^2}
\]
两边都是非负数，所以平方不会引入增根，从而
\[
(x-1)^2+(y-3)^2=(x+3)^2+(y-5)^2
\]
展开并消去两边的 \(x^2\) 和 \(y^2\)，得
\[
x^2-2x+1+y^2-6y+9=x^2+6x+9+y^2-10y+25
\]
即 \(-8x+4y-24=0\)，化简为 \(2x-y+6=0\).

反过来，任意满足 \(2x-y+6=0\) 的点都满足上述两个距离平方相等；由于距离非负，
两个距离也相等. 因此所求动点的轨迹方程为 \(2x-y+6=0\).
\end{solution}
